\chapter{Numbers, Characters and Symbols}\label{ch:numbers}

This chapter describes Dylan's simplest objects: numbers, characters,
symbols and keywords. It follows chapters 14, 15 and 16 of the 1992
manual (\dylanbook{chapters 14--16}) and, for comparison, chapter~9
(\dylanbook{chapter 9}). The two implementations agree on almost
everything here, but numbers are where they differ most: the Common Lisp
implementation has the unbounded integers and the complex numbers the
manual describes, and the compiler to C has neither. Each difference is
pointed out where it arises; appendix~\ref{app:differences} collects them.

The examples show their results as comments, \code{; => value}. Where the
two implementations print a result differently, both are shown.

\section{The number classes}\label{sec:numbers-classes}

\index{number classes}
Numbers are arranged in a hierarchy of classes that follows the
mathematical kinds of number. The abstract classes have no direct
instances; they exist so that methods can specialize on them.

\begin{center}
\begin{tabular}{lll}
\toprule
Class & Direct superclass & Kind \\
\midrule
\code{<number>}\index{<number>} & \code{<object>} & abstract \\
\code{<complex>}\index{<complex>} & \code{<number>} & abstract \\
\code{<real>}\index{<real>} & \code{<complex>} & abstract \\
\code{<rational>}\index{<rational>} & \code{<real>} & abstract \\
\code{<integer>}\index{<integer>} & \code{<rational>} & concrete \\
\code{<ratio>}\index{<ratio>} & \code{<rational>} & concrete \\
\code{<float>}\index{<float>} & \code{<real>} & abstract \\
\code{<single-float>}\index{<single-float>} & \code{<float>} & concrete \\
\code{<double-float>}\index{<double-float>} & \code{<float>} & concrete \\
\code{<extended-float>}\index{<extended-float>} & \code{<float>} & concrete \\
\code{<rectangular-complex>}\index{<rectangular-complex>} & \code{<complex>} & concrete; Common Lisp only \\
\bottomrule
\end{tabular}
\end{center}

\noindent All of these classes but \code{<number>} are sealed: a program
may not define new subclasses of them, though it may define methods on
them. Every number
belongs to exactly one of the concrete classes, which
\code{object-class} reports:

\begin{dylan}
(object-class 3)        ; => <integer>
(object-class 1/2)      ; => <ratio>
(object-class (/ 4 2))  ; => <integer>
(object-class 1.5)      ; => <double-float>
(object-class 1.5s0)    ; => <single-float>
\end{dylan}

\noindent (A class prints as \code{\{the class <integer>\}}; this
chapter writes just the name.)

\subsection{Integers}

\index{integer!range}\index{overflow}
An \code{<integer>} is a whole number. The two implementations differ in
how large one may be:

\begin{itemize}
\item In the \textbf{Common Lisp implementation}, integers are unbounded,
  as the manual requires: \code{(expt 2 100)} is
  \code{1267650600228229401496703205376}, and an integer of any size is
  still an \code{<integer>}.
\item In the \textbf{compiler to C}, an integer is 63 bits, from
  $-2^{62}$ to $2^{62}-1$, that is, from \code{-4611686018427387904} to
  \code{4611686018427387903}. Every integer operation checks for
  overflow, and a result outside that range signals an error rather than
  giving a wrong answer:
\end{itemize}

\begin{dylan}
(+ 4611686018427387903 1)
; C:    Dylan error: integer overflow (integers are 63 bits)
; Lisp: => 4611686018427387904
\end{dylan}

\noindent The error is an ordinary \code{<simple-error>} (chapter~\ref{ch:conditions}),
so a program can handle it:

\begin{dylan}
(handler-case (* 4611686018427387903 2)
  ((<error>) 'overflow))
; C:    => overflow
; Lisp: => 9223372036854775806
\end{dylan}

\noindent The check applies to every operation that makes an integer:
\code{+}, \code{-}, \code{*}, \code{expt}, \code{ash}, \code{lcm}, and
the numerators and denominators of ratios.

\subsection{Ratios}

\index{ratio}
A \code{<ratio>} is an exact fraction, kept in lowest terms with a
positive denominator. Dividing integers that do not divide evenly makes
one, and a ratio whose denominator would be 1 is made an integer instead
(\term{rational canonicalization}):

\begin{dylan}
(/ 1 3)         ; => 1/3
(/ 6 3)         ; => 2
(/ -6 4)        ; => -3/2
(+ 1/3 1/6)     ; => 1/2
(+ 1/2 1/2)     ; => 1
(* 1/2 4)       ; => 2
\end{dylan}

\noindent In the compiler to C the numerator and the denominator are
each limited to 63 bits, with the same overflow error as integers;
intermediate results are computed exactly in 128 bits, so only a final
result that does not fit is an error. In the Common Lisp implementation
ratios are unbounded.

\subsection{Floats}

\index{float}
The three float classes are the IEEE formats the manual names. Their
representations are:

\begin{center}
\begin{tabular}{lll}
\toprule
Class & Compiler to C & Common Lisp implementation \\
\midrule
\code{<single-float>} & C \code{float} (32 bits) & SBCL \code{single-float} \\
\code{<double-float>} & C \code{double} (64 bits) & SBCL \code{double-float} \\
\code{<extended-float>} & C \code{double}, a class of its own & no instances \\
\bottomrule
\end{tabular}
\end{center}

\noindent \code{<double-float>} is the default: a float literal without
an exponent marker, or with \code{e} or \code{d}, is a double. In the
compiler to C an \code{<extended-float>} is a separate class with a
double's precision; the only way to make one is \code{as} (see
section~\ref{sec:numbers-conversion}). In the Common Lisp implementation,
SBCL's long float is its double float, so \code{<extended-float>} has no
instances at all, and \code{(as <extended-float> 1)} answers the
\code{<double-float>} \code{1.0}.

\subsection{Complex numbers}

\index{complex numbers}
The manual's \code{<complex>} numbers exist only in the Common Lisp
implementation, where they are Common Lisp's complex numbers and belong
to the class \code{<rectangular-complex>}:

\begin{dylan}
(make-rectangular 1 2)                   ; Lisp: => #c(1 2)
(make <complex> real: 1 imag: 2)         ; Lisp: => #c(1 2)
(* (make-rectangular 0 1)
   (make-rectangular 0 1))               ; Lisp: => -1
(abs (make-rectangular 3 4))             ; Lisp: => 5.0
(sqrt -4)                                ; Lisp: => #c(0.0 2.0)
\end{dylan}

\noindent The init-keywords of \code{<complex>} are \code{real:},
\code{imag:}, \code{magnitude:} and \code{angle:}, as in
\dylanbook{chapter 14}; \code{make-rectangular} and \code{make-polar}
make complex numbers from parts, and \code{real-part},
\code{imag-part} and \code{angle} take them apart. As the manual
requires, a complex number whose imaginary part is an exact zero is
canonicalized to a real: \code{(make-rectangular 3 0)} is \code{3}.

In the compiler to C, \code{<complex>} is kept, abstract, as the
superclass of \code{<real>}, but it has no instances of its own.
\code{make-rectangular} and \code{make-polar} accept only an exact zero
for the imaginary part or the angle, and answer the real; anything else
is an error, as is \code{make} of \code{<complex>} and the square root
of a negative number:

\begin{dylan}
(make-rectangular 3 0)   ; => 3
(make-rectangular 1 2)   ; C: error: complex numbers are not supported: 1 2
(sqrt -4)                ; C: error: sqrt: the result would be complex: -4
\end{dylan}

\section{Writing numbers}\label{sec:numbers-literals}

\index{number literals}
Numbers are written as in Common Lisp:

\begin{center}
\begin{tabular}{lll}
\toprule
Written & Value & Class \\
\midrule
\code{42}, \code{-7} & 42, $-7$ & \code{<integer>} \\
\code{1.} & 1 & \code{<integer>} (a trailing point means decimal) \\
\code{\#xFF}, \code{\#b101}, \code{\#o17} & 255, 5, 15 & \code{<integer>} \\
\code{6/4} & \code{3/2} & \code{<ratio>} \\
\code{1.5}, \code{.5}, \code{-0.0} & & \code{<double-float>} \\
\code{1e3}, \code{2d0} & \code{1000.0}, \code{2.0} & \code{<double-float>} \\
\code{1.5s2}, \code{1.5f2} & \code{150.0s0} & \code{<single-float>} \\
\code{1.5l0} & \code{1.5} & \code{<double-float>} \\
\bottomrule
\end{tabular}
\end{center}

\noindent The exponent markers \code{e} and \code{d} make a double,
\code{s} and \code{f} a single. The marker \code{l} also makes a double,
in both implementations; there is no way to write an
\code{<extended-float>} literal. A ratio is written with no spaces,
\code{numerator/denominator}, and is reduced as it is read.

\section{Arithmetic}\label{sec:numbers-arithmetic}

\index{arithmetic}\index{+}\index{-}\index{*}\index{/}
The four arithmetic functions take any number of arguments, as in
\dylanbook{chapter 14}:

\begin{dylan}
(+)             ; => 0
(*)             ; => 1
(+ 1 2 3)       ; => 6
(- 5)           ; => -5         the additive inverse
(- 10 1 2 3)    ; => 4          10 - 1 - 2 - 3
(/ 2)           ; => 1/2        the multiplicative inverse
(/ 0.5)         ; => 2.0
(/ 60 2 3)      ; => 10         60 / 2 / 3
(/ 7 2)         ; => 7/2
(/ 7.0 2)       ; => 3.5
\end{dylan}

\noindent \code{-} and \code{/} need at least one argument. With several
arguments, each function combines them from left to right by calling
\term{binary+}, \code{binary-}\index{binary-}, \code{binary*}\index{binary*}
or \code{binary/}\index{binary/}; with one, \code{-} calls
\code{unary-}\index{unary-} and \code{/} calls \code{unary/}\index{unary/}.
The variadic functions themselves are methods and cannot be extended,
but the binary and unary ones are generic functions: to make
\code{+} work on a class of your own, add a method to \code{binary+}.

\begin{dylan}
(define-class <money> (<object>)
  (cents init-keyword: cents:))

(define-method binary+ ((a <money>) (b <money>))
  (make <money> cents: (+ (cents a) (cents b))))

(cents (+ (make <money> cents: 150)
          (make <money> cents: 275)
          (make <money> cents: 5)))      ; => 430
\end{dylan}

\index{division by zero}
Dividing by an exact or a floating zero is an error in both
implementations, as design note~24 requires; so is \code{floor/},
\code{modulo} or any of the division family with a zero divisor:

\begin{dylan}
(handler-case (/ 1 0)   ((<error>) 'caught))   ; => caught
(handler-case (/ 1.0 0) ((<error>) 'caught))   ; => caught
\end{dylan}

\subsection{Contagion}

\index{contagion}
When the arguments of an arithmetic function are of different classes,
the rules are Common Lisp's (\dylanbook{chapter 14}): an operation
between a rational and a float answers a float of the float's format,
and an operation between two floats answers the wider format (single is
narrower than double, double than extended). Rational results are
canonicalized. Because a call with several arguments is a chain of
binary calls, the conversions happen from left to right.

\begin{dylan}
(+ 1 2.0)          ; => 3.0
(+ 1/2 0.5)        ; => 1.0
(* 2 0.5s0)        ; => 1.0s0 (C)   1.0f0 (Lisp)
(+ 1.0s0 1.0d0)    ; => 2.0         a <double-float>
(* 1.5s0 1.5d0)    ; => 2.25        a <double-float>
\end{dylan}

\begin{notebox}
The compiler to C converts a rational to the float's format by way of a
C \code{double}, so a single-float result can differ from the Common Lisp
implementation's in the last place: \code{(+ 1/3 0.5s0)} is
\code{0.8333333s0} in compiled code and \code{0.8333334f0} in the Common
Lisp implementation.
\end{notebox}

\section{Comparing numbers}\label{sec:numbers-comparison}

\index{comparison}\index{=}\index{<}
The comparison functions of \dylanbook{chapter 9} take two or more
arguments and compare them pairwise from left to right:

\begin{center}
\begin{tabular}{ll}
\toprule
Function & True when \\
\midrule
\code{(= a b c ...)} & every argument is \code{=} to the next \\
\code{(/= a b)} & the two arguments are not \code{=} (exactly two) \\
\code{(< a b c ...)} & each argument is less than the next \\
\code{(> a b c ...)} & each argument is greater than the next \\
\code{(<= a b c ...)} & each argument is less than or equal to the next \\
\code{(>= a b c ...)} & each argument is greater than or equal to the next \\
\bottomrule
\end{tabular}
\end{center}

\begin{dylan}
(< 1 2 3)       ; => #t
(< 1 3 2)       ; => #f
(<= 1 1 2)      ; => #t
(>= 3 3 1)      ; => #t
(= 3 3 3)       ; => #t
(/= 1 2)        ; => #t
\end{dylan}

\noindent Numbers are compared by mathematical value, whatever their
classes. A float is compared with a rational exactly (\term{rational
contagion for comparisons}): the float's exact value is used, not a
rounded one. Since the double nearest $0.1$ is not exactly $1/10$, they
are not \code{=}:

\begin{dylan}
(= 1 1.0)                     ; => #t
(= 1/2 0.5)                   ; => #t
(= 0.1 1/10)                  ; => #f
(< 1/3 0.3333333333333333)    ; => #f
\end{dylan}

\index{id?}
\code{=} is not \code{id?}. Two numbers of the same class and value are
\code{id?}; numbers of different classes never are, even when they are
\code{=}:

\begin{dylan}
(id? 1 1)       ; => #t
(id? 1 1.0)     ; => #f
\end{dylan}

\noindent The comparison functions are methods; to extend them to a
class of your own, add methods to the generic functions
\code{binary=}\index{binary=} and \code{binary<}\index{binary<}, on which
all six are built (section~\ref{sec:collections-equality} shows an
example).

\section{Predicates}\label{sec:numbers-predicates}

\index{predicates!numeric}
These generic functions test a number and answer \code{\#t} or
\code{\#f}:

\begin{center}
\begin{tabular}{lll}
\toprule
Function & Argument & True when \\
\midrule
\code{zero?}\index{zero?} & a number & it is zero (\code{0}, \code{0.0}, \code{-0.0}) \\
\code{positive?}\index{positive?} & a real & it is greater than zero \\
\code{negative?}\index{negative?} & a real & it is less than zero \\
\code{odd?}\index{odd?} & an integer & it is odd \\
\code{even?}\index{even?} & an integer & it is even \\
\code{integral?}\index{integral?} & a number & it has an integer value \\
\bottomrule
\end{tabular}
\end{center}

\begin{dylan}
(zero? 0.0)         ; => #t
(positive? -1)      ; => #f
(negative? -1/2)    ; => #t
(odd? 3)            ; => #t
(even? 3)           ; => #f
(integral? 3)       ; => #t
(integral? 2.0)     ; => #t
(integral? 2.5)     ; => #f
(integral? 1/2)     ; => #f
\end{dylan}

\noindent To ask which kind of number an object is, use \code{instance?}
with the classes of section~\ref{sec:numbers-classes}:

\begin{dylan}
(instance? 3 <rational>)     ; => #t
(instance? 1/2 <rational>)   ; => #t
(instance? 1.5 <real>)       ; => #t
(instance? 1.5 <integer>)    ; => #f
\end{dylan}

\section{The division family}\label{sec:numbers-division}

\index{division}
Dylan has the rounding and division functions of Common Lisp.
Each of the one-argument functions rounds a real to an integer and
answers two values: the integer and the remainder, which is the
difference between the argument and the integer. They differ in which
way they round:

\begin{center}
\begin{tabular}{ll}
\toprule
Function & Rounds toward \\
\midrule
\code{floor}\index{floor} & negative infinity \\
\code{ceiling}\index{ceiling} & positive infinity \\
\code{truncate}\index{truncate} & zero \\
\code{round}\index{round} & the nearest integer; a tie goes to the even one \\
\bottomrule
\end{tabular}
\end{center}

\begin{dylan}
(floor 7/2)      ; => 3 and 1/2
(floor -7/2)     ; => -4 and 1/2
(ceiling 2.5)    ; => 3 and -0.5
(truncate -2.7)  ; => -2 and -0.7000000000000002
(round 2.5)      ; => 2 and 0.5       the tie goes to 2, which is even
(round 3.5)      ; => 4 and -0.5
(round -2.5)     ; => -2 and -0.5
\end{dylan}

\noindent Use \code{bind} to receive both values (chapter~\ref{ch:control});
where one value is wanted, the first is used:

\begin{dylan}
(bind ((q r (floor 7/2)))
  (list q r))                ; => (3 1/2)
(+ (round 7/2) 1)            ; => 5
\end{dylan}

\index{floor/}\index{ceiling/}\index{round/}\index{truncate/}
The two-argument functions \code{floor/}, \code{ceiling/},
\code{round/} and \code{truncate/} divide their first argument by their
second, round the quotient as the one-argument function of the same name
does, and answer the rounded quotient and the remainder:

\begin{dylan}
(floor/ 7 2)        ; => 3 and 1
(floor/ -7 2)       ; => -4 and 1
(ceiling/ 7 2)      ; => 4 and -1
(truncate/ -7 2)    ; => -3 and -1
(round/ 7 2)        ; => 4 and -1
(round/ 5 2)        ; => 2 and 1
(floor/ 7.5 2)      ; => 3 and 1.5
\end{dylan}

\index{modulo}\index{remainder}
\code{modulo} answers the remainder of \code{floor/}, which has the sign
of the divisor; \code{remainder} answers the remainder of
\code{truncate/}, which has the sign of the dividend:

\begin{dylan}
(modulo 7 2)        ; => 1
(modulo -7 2)       ; => 1
(modulo 7 -2)       ; => -1
(remainder 7 2)     ; => 1
(remainder -7 2)    ; => -1
(remainder 7 -2)    ; => 1
\end{dylan}

\section{gcd, lcm, abs, min and max}\label{sec:numbers-misc}

\index{gcd}\index{lcm}
\code{gcd} answers the greatest common divisor of its integer arguments
and \code{lcm} their least common multiple. With no arguments they
answer \code{0} and \code{1}; with one, the argument. They are built on
the generic functions \code{binary-gcd}\index{binary-gcd} and
\code{binary-lcm}\index{binary-lcm}, which programs can extend.

\begin{dylan}
(gcd)            ; => 0
(gcd 12 18)      ; => 6
(gcd 12 18 8)    ; => 2
(gcd -12 18)     ; => 6
(lcm)            ; => 1
(lcm 4 6)        ; => 12
(lcm 4 6 10)     ; => 60
(lcm 0 5)        ; => 0
\end{dylan}

\index{abs}
\code{abs} answers the absolute value of a number (for a complex number,
in the Common Lisp implementation, its magnitude); \code{magnitude}\index{magnitude}
is the same for a real.

\begin{dylan}
(abs -5)       ; => 5
(abs -5/3)     ; => 5/3
(abs -2.5)     ; => 2.5
(magnitude -3) ; => 3
\end{dylan}

\index{min}\index{max}
\code{min} and \code{max} take one or more reals and answer the least or
the greatest of them, comparing with \code{binary<}. The answer is one
of the arguments, not converted:

\begin{dylan}
(min 3 1 2)       ; => 1
(max 3 1 2)       ; => 3
(max 1/2 0.4)     ; => 1/2
(max 1 2.0)       ; => 2.0
(min 3 1.5s0)     ; => 1.5s0 (C)   1.5f0 (Lisp)
(min 2)           ; => 2
\end{dylan}

\noindent When several arguments are equal the two implementations may
choose different ones: \code{(min 1 1.0)} is \code{1} in the compiler to
C, which keeps the first of equal arguments, and \code{1.0} in the
Common Lisp implementation, which uses Common Lisp's \code{min}.
Because \code{min} and \code{max} use \code{binary<}, they work on any
objects \code{<} can compare, characters among them:
\code{(max \#\textbackslash{}a \#\textbackslash{}z)} is \code{\#\textbackslash{}z}.

\section{Exponentials and trigonometry}\label{sec:numbers-transcendental}

\index{sqrt}\index{exp}\index{log}\index{expt}
\index{trigonometric functions}
The irrational functions of \dylanbook{chapter 14} are defined as in
Common Lisp:

\begin{center}
\begin{tabular}{ll}
\toprule
Function & Answers \\
\midrule
\code{(sqrt x)} & the square root of $x$ \\
\code{(exp x)} & $e^x$ \\
\code{(log x)} & the natural logarithm of $x$ \\
\code{(expt x y)} & $x^y$ \\
\code{(sin x)}, \code{(cos x)}, \code{(tan x)} & the trigonometric functions, $x$ in radians \\
\code{(asin x)}, \code{(acos x)}, \code{(atan x)} & their inverses \\
\code{(atan2 y x)} & the angle of the point $(x, y)$, from $-\pi$ to $\pi$ \\
\code{(sinh x)}, \code{(cosh x)}, \code{(tanh x)} & the hyperbolic functions \\
\code{(asinh x)}, \code{(acosh x)}, \code{(atanh x)} & their inverses \\
\bottomrule
\end{tabular}
\end{center}

\noindent For a rational argument they answer a \code{<double-float>}; for
a float argument, a float of the same format. So \code{(sqrt 25)} is
\code{5.0}, not \code{5}: the manual's examples show these functions
answering floats, and so do both implementations.

\begin{dylan}
(sqrt 25)        ; => 5.0
(sqrt 2)         ; => 1.4142135623730951
(sqrt 2.0s0)     ; => 1.4142135s0 (C)   1.4142135f0 (Lisp)
(exp 1)          ; => 2.718281828459045
(log 1)          ; => 0.0
(log 100)        ; => 4.605170185988092
(sin 0)          ; => 0.0
(cos 0)          ; => 1.0
(atan 1)         ; => 0.7853981633974483
(atan2 1 1)      ; => 0.7853981633974483
(atan2 -1 -1)    ; => -2.356194490192345
(asin 1)         ; => 1.5707963267948966
(acosh 1)        ; => 0.0
\end{dylan}

\noindent \code{expt} is exact when its base is rational and its power
is an integer, and answers a float otherwise:

\begin{dylan}
(expt 2 10)      ; => 1024
(expt 2 -2)      ; => 1/4
(expt 2/3 3)     ; => 8/27
(expt 10 0)      ; => 1
(expt 2.0 3)     ; => 8.0
(expt 4 1/2)     ; => 2.0
(expt 2.0 0.5)   ; => 1.4142135623730951
\end{dylan}

\noindent Where the true result is not a real number, or not finite, the
implementations behave differently:

\begin{center}
\begin{tabular}{lp{0.25\textwidth}p{0.43\textwidth}}
\toprule
Expression & Compiler to C & Common Lisp implementation \\
\midrule
\code{(sqrt -4)} & error: the result would be complex & \code{\#c(0.0 2.0)} \\
\code{(log 0)} & error: the argument must be positive & error (Lisp's division by zero) \\
\code{(exp 1000)} & \code{+infinity} & error (Lisp's floating-point overflow) \\
\code{(expt 2 100)} & error: integer overflow & \code{1267650600228229401496703205376} \\
\bottomrule
\end{tabular}
\end{center}

\noindent An error from Common Lisp's arithmetic reaches a Dylan
handler as a \code{<simple-error>} whose message begins
\code{Lisp condition:} (chapter~\ref{ch:conditions}).

\section{Integers as bits}\label{sec:numbers-logical}

\index{logical operations}
These functions treat an integer as an infinite two's-complement string
of bits, as Common Lisp does. \code{logior}, \code{logxor},
\code{logand} and \code{logeqv} take any number of integers; with none,
they answer their identities, \code{0}, \code{0}, \code{-1} and
\code{-1}. The rest take exactly the arguments shown.

\begin{center}
\begin{tabular}{ll}
\toprule
Function & Answers \\
\midrule
\code{(logior i ...)}\index{logior} & inclusive or \\
\code{(logxor i ...)}\index{logxor} & exclusive or \\
\code{(logand i ...)}\index{logand} & and \\
\code{(logeqv i ...)}\index{logeqv} & equivalence (exclusive nor) \\
\code{(lognot i)}\index{lognot} & the complement \\
\code{(lognand i j)}\index{lognand} & \code{(lognot (logand i j))} \\
\code{(lognor i j)}\index{lognor} & \code{(lognot (logior i j))} \\
\code{(logandc1 i j)}\index{logandc1} & \code{(logand (lognot i) j)} \\
\code{(logandc2 i j)}\index{logandc2} & \code{(logand i (lognot j))} \\
\code{(logorc1 i j)}\index{logorc1} & \code{(logior (lognot i) j)} \\
\code{(logorc2 i j)}\index{logorc2} & \code{(logior i (lognot j))} \\
\code{(logbit? index i)}\index{logbit?} & whether bit \meta{index} of \meta{i} is 1 \\
\code{(ash i count)}\index{ash} & \meta{i} shifted left \meta{count} bits (right if negative) \\
\bottomrule
\end{tabular}
\end{center}

\begin{dylan}
(logior 1 2 4)      ; => 7
(logand 12 10)      ; => 8
(logxor 12 10)      ; => 6
(logeqv 12 10)      ; => -7
(lognot 5)          ; => -6
(lognand 12 10)     ; => -9
(lognor 12 10)      ; => -15
(logandc1 12 10)    ; => 2
(logandc2 12 10)    ; => 4
(logorc1 12 10)     ; => -5
(logorc2 12 10)     ; => -3
(logbit? 0 5)       ; => #t
(logbit? 1 5)       ; => #f
(ash 1 10)          ; => 1024
(ash 1024 -3)       ; => 128
(ash -8 -1)         ; => -4
\end{dylan}

\noindent In the compiler to C a shift whose result does not fit in 63
bits is an overflow error: \code{(ash 1 70)} is an error there and
\code{1180591620717411303424} in the Common Lisp implementation.

\section{Parts of numbers}\label{sec:numbers-parts}

\index{numerator}\index{denominator}\index{rationalize}
\code{numerator} and \code{denominator} answer the parts of a rational
in lowest terms; an integer's denominator is 1. For a float, both
implementations answer the parts of the fraction the float holds exactly,
as floats:

\begin{dylan}
(numerator 6/4)       ; => 3
(denominator 6/4)     ; => 2
(denominator 5)       ; => 1
(numerator 1.5)       ; => 3.0
(denominator 1.5)     ; => 2.0
(denominator 0.1)     ; => 3.602879701896397e16
\end{dylan}

\noindent The last line shows that the double nearest $0.1$ is a fraction
with a large power of two below it. \code{rationalize} answers the
simplest rational that is within the float's precision of it, which is
usually the fraction you meant:

\begin{dylan}
(rationalize 0.5)     ; => 1/2
(rationalize 0.1)     ; => 1/10
(rationalize 0.333)   ; => 333/1000
(rationalize 1/3)     ; => 1/3
\end{dylan}

\index{real-part}\index{imag-part}\index{angle}
\code{real-part}, \code{imag-part} and \code{angle} work on reals in
both implementations (and on complex numbers in the Common Lisp
implementation): a real is its own real part, its imaginary part is
\code{0}, and its angle is zero for a non-negative number and $\pi$ for a
negative one.

\begin{dylan}
(real-part 5)    ; => 5
(imag-part 5)    ; => 0
(angle 5)        ; => 0 (C)   0.0 (Lisp)
(angle -5)       ; => 3.141592653589793
\end{dylan}

\section{Converting numbers}\label{sec:numbers-conversion}

\index{as!numbers}
\code{as} converts a real to a float of a given format:

\begin{dylan}
(as <float> 1/2)           ; => 0.5        a <double-float>
(as <double-float> 1)      ; => 1.0
(as <single-float> 1/3)    ; => 0.33333334s0 (C)   0.33333334f0 (Lisp)
(as <extended-float> 1)    ; => 1.0l0 (C)   1.0, a <double-float> (Lisp)
\end{dylan}

\noindent To go from a float to an integer, round it with \code{floor},
\code{ceiling}, \code{round} or \code{truncate}; to go to a rational,
use \code{rationalize}, or \code{numerator} and \code{denominator}. The
compiler to C accepts \code{as} only for these float conversions (and
for a number already of the class asked for): \code{(as <integer> 3.0)}
and \code{(as <rational> 0.5)} are errors there. The Common Lisp
implementation also converts a float with an integer value to an
integer, \code{(as <integer> 3.0)} answering \code{3}, and any float to
its exact rational, \code{(as <rational> 0.5)} answering \code{1/2}; it
signals an error for \code{(as <integer> 2.5)}.

\section{How floats print}\label{sec:numbers-printing}

\index{float!printing}
Both implementations print a float with the fewest digits that read back
as the same float, so a float printed and read again is the same float.
A float prints in fixed notation when its magnitude is at least
$10^{-5}$ and less than $10^{16}$, or it is zero, and otherwise with an
exponent. A double-float, the default format, is marked only by its
exponent, written with \code{e}; a single-float ends in \code{s0} (or
\code{s}\meta{n} with an exponent), an extended-float in \code{l0}. An
infinity prints as \code{+infinity} or \code{-infinity} and a NaN as
\code{+nan}; neither can be read back.

\begin{center}
\begin{tabular}{ll}
\toprule
Value & Prints as \\
\midrule
\code{0.1} & \code{0.1} \\
\code{(/ 1.0 3)} & \code{0.3333333333333333} \\
\code{100.0} & \code{100.0} \\
\code{123456789.0} & \code{123456789.0} \\
\code{1e15} & \code{1000000000000000.0} \\
\code{1e16} & \code{1.0e16} \\
\code{1e-4} & \code{0.0001} \\
\code{1e-7} & \code{1.0e-7} \\
\code{(/ 2.0s0 3)} & \code{0.6666667s0} \\
\code{(* 1.5s0 2)} & \code{3.0s0} \\
\bottomrule
\end{tabular}
\end{center}

\noindent Integers print in decimal and ratios as
\code{numerator/denominator}. Printing in general is
the subject of chapter~\ref{ch:printing}.

\section{Characters}\label{sec:characters}

\index{<character>}\index{character}
A \code{<character>} is written \code{\#\textbackslash{}} followed by the
character, or by one of the names \code{space}, \code{newline},
\code{tab}, \code{page}, \code{rubout} and \code{backspace}:

\begin{dylan}
(list #\a #\A #\( #\space #\newline #\tab)
; => (#\a #\A #\( #\space #\newline #\tab)
\end{dylan}

\index{as!characters}
A character's integer is its Unicode code point, in both
implementations. \code{as} converts in either direction
(\dylanbook{chapter 16}):

\begin{dylan}
(as <integer> #\A)        ; => 65
(as <integer> #\newline)  ; => 10
(as <character> 97)       ; => #\a
\end{dylan}

\index{as-uppercase}\index{as-lowercase}
\code{as-uppercase} and \code{as-lowercase} answer the other case of a
character, or the character itself if it has no other case. Both
implementations follow Unicode, not only ASCII: the uppercase of Greek
small lambda, code point 955, is capital lambda, 923.

\begin{dylan}
(as-uppercase #\a)    ; => #\A
(as-lowercase #\Z)    ; => #\z
(as-uppercase #\1)    ; => #\1
(as <integer> (as-uppercase (as <character> 955)))   ; => 923
\end{dylan}

\index{character!comparison}
Characters compare by code point, so case is significant and every
capital ASCII letter is less than every small one:

\begin{dylan}
(< #\a #\b)       ; => #t
(< #\A #\a)       ; => #t
(> #\a #\B)       ; => #t
(= #\a #\A)       ; => #f
(id? #\a #\a)     ; => #t
\end{dylan}

\noindent \code{as-uppercase} and \code{as-lowercase} also work on whole
strings (section~\ref{sec:collections-strings}), which is how to compare
strings without regard to case. A string of characters that do not all
fit in eight bits is a \code{<unicode-string>} in the compiler to C; in
the Common Lisp implementation every string is a \code{<byte-string>},
which holds any character.

\section{Symbols and keywords}\label{sec:symbols}

\index{<symbol>}\index{<keyword>}\index{symbol}\index{keyword}
Symbols and keywords are objects identified by their names
(\dylanbook{chapter 15}). A symbol is written as a name, quoted when it
is to be a value rather than a variable: \code{'apple}. A keyword is a
name followed by a colon, \code{apple:}, and evaluates to itself; the
colon is not part of its name. Symbols name variables; keywords name
keyword arguments and slot options.

\begin{dylan}
(object-class 'apple)     ; => <symbol>
(object-class apple:)     ; => <keyword>
\end{dylan}

\index{symbol!case}
Names are not case sensitive: \code{Foo}, \code{foo} and \code{FOO} are
one symbol, and two symbols with the same name are always \code{id?}.
A symbol does remember a spelling, the one it was first written with,
and prints that way. Converting between symbols, keywords and strings is
done with \code{as}:

\begin{dylan}
(list 'Hello 'hello 'HELLO)           ; => (Hello Hello Hello)
(id? 'Hello 'HELLO)                   ; => #t
(as <symbol> "hello")                 ; => Hello
(as <string> 'hello)                  ; => "Hello"
(as <keyword> "colour")               ; => colour:
(as <string> 'colour:)                ; => "colour"
(as <symbol> "two words")             ; => two words
\end{dylan}

\noindent The spelling remembered is the first anywhere, and the
implementation's own names come first. In the compiler to C that
includes every name the library uses, so \code{'Size} prints as
\code{size} and \code{Test:} as \code{test:} there, while the Common
Lisp implementation, which predefines fewer names, prints them
\code{Size} and \code{Test:}. Either way \code{(id? 'size 'Size)} is
\code{\#t}; only the printing differs.

\index{symbol!comparison}
Symbols are compared only by identity: \code{(= 'abc 'ABC)} and
\code{(id? 'abc 'ABC)} are both \code{\#t}, and there is no ordering, so
\code{(< 'a 'b)} is an error (no applicable method for
\code{binary<}). To sort symbols, sort their names:

\begin{dylan}
(sort '(pear apple fig)
      test: (method (a b) (< (as <string> a) (as <string> b))))
; => (apple fig pear)
\end{dylan}
